\medskip

L'espace est muni d'un repère orthonormé \Oijk.

On considère les points

\begin{center}A$\left(2\sqrt 3~;~0~;~0\right)$,\quad  B(0~;~2~;~0), \quad C(0~;~0~;~1)\quad et \quad K$\left(\dfrac{\sqrt 3}{2}~;~\dfrac32~;~0\right)$.\end{center}

\psset{unit=1cm,arrowsize=2pt 3}
\begin{center}
\begin{pspicture}(-4,-2)(4,2.3)

\psline(-4,-2)\psline(3.9,-0.7)\psline(0,2.3)
\psline[linecolor=blue]{->}(-1.1,-0.55)\psline[linecolor=blue]{->}(1.7,-0.285)\psline[linecolor=blue]{->}(0,1.8)
\pspolygon(-3.5,-1.73)(3.4,-0.6)(0,1.8)
\psline(0,1.8)(1.55,-0.9)
\uput[ul](-3.5,-1.73){A} \uput[ur](3.4,-0.6){B} \uput[ul](0,1.8){C} \uput[d](0,0){O} \uput[dr](1.55,-0.9){K}
\uput[u](-0.5,-0.3){\blue $\vect{\imath}$} \uput[u](0.6,-0.1){\blue $\vect{\jmath}$} \uput[l](0,0.8){\blue $\vect{k}$}
\end{pspicture}
\end{center}

\medskip

\begin{enumerate}
\item Justifier qu'une représentation paramétrique de la droite (CK) est :

\renewcommand\arraystretch{1.7}
\[\left\{\begin{array}{l c l}
x&=&\phantom{1-}\dfrac{\sqrt 3}{2}t\\
y&=&\phantom{1-}\dfrac32 t\\
z&=&1 - t
\end{array}\right. \: (t \in \R)\]
\renewcommand\arraystretch{1.}

\item Soit M$(t)$ un point de la droite (CK) paramétrée par un réel $t$.

Établir que OM$(t) = \sqrt{4t^2 - 2t + 1}$.
\item Soit $f$ la fonction définie et dérivable sur $\R$ par $f(t) = \text{OM}(t)$.
	\begin{enumerate}
		\item Étudier les variations de la fonction $f$ sur $\R$.
		\item En déduire la valeur de $t$ pour laquelle $f$ atteint son minimum.
	\end{enumerate}
\item En déduire que le point H$\left(\dfrac{\sqrt3}{8}~;~\dfrac38~;~\dfrac34\right)$ est le projeté orthogonal du point O sur la droite~(CK).
\item Démontrer, à l'aide de l'outil produit scalaire, que le point H est l'orthocentre (intersection des hauteurs d'un triangle) du triangle ABC.
\item 
	\begin{enumerate}
		\item Démontrer que la droite (OH) est orthogonale au plan (ABC).
		\item En déduire une équation du plan (ABC). 
	\end{enumerate}
	\item Calculer, en unité d'aire, l'aire du triangle ABC.
\end{enumerate}
